\chapter{我们其实是一个生态系统}

\begin{kaichang}
你家冰箱里多半有一杯酸奶。拿起来看看配料表：生牛乳、白砂糖，还有一行小字——保加利亚乳杆菌、嗜热链球菌。注意，这不是防腐剂，不是香精，而是\textbf{活的细菌}。厂家故意把细菌放进牛奶里，你还花钱买，喝完还觉得挺健康。

现在做一件更需要勇气的事：张开嘴，对着镜子看看自己的舌头。那上面此刻住着几百种微生物，数量以亿计。再往下想一层：你的肠道里，住着一个以万亿计的微生物世界。有人甚至说，你身体里的细菌细胞，数量和你自己的细胞差不多。

先猜一猜：如果把你全身上下所有的细菌收集起来称重，你觉得有多重？一千克？十千克？还是只有一小勺？这一章结束时，我们回来对答案——顺便回答一个更大的问题：你究竟是一个“人”，还是一个会走路的生态系统？
\end{kaichang}

\section{一杯酸奶里的活宇宙}

先做现象层该做的事：只描述看得见、闻得到、测得出的东西。

新鲜牛奶是甜的，有淡淡的奶香，放久了会坏掉发臭。但如果在合适的温度下接入合适的细菌，牛奶会在几个小时内变稠、变酸，变成酸奶。酸味来自乳酸——细菌吃下牛奶里的糖，吐出乳酸，酸又让牛奶里的蛋白质凝固（蛋白质遇酸变性，第 47 章讲过这件事）。你尝到的每一口酸、感受到的每一丝浓稠，都是细菌代谢的“排泄物”和它的副产品。

再看你的身体。你的皮肤表面摸起来干干净净，但用显微镜看，它像一片起伏的平原，沟壑里住着葡萄球菌、丙酸杆菌；你的口腔湿润温暖，是几百种细菌的“热带雨林”；而你的大肠——黑暗、温暖、缺氧、食物残渣源源不断地流过——是整个身体里最繁华的微生物都市。

其实它们也留下了不少“可观测痕迹”，只是你没往那方面想。早晨醒来嘴里那股味道，是口腔细菌连夜分解食物残渣的产物的气味；汗水本身几乎没有味道，腋下的气味来自皮肤细菌分解汗液成分后放出的小分子；刷牙时从牙刷上刮下来的那层软垢——牙菌斑——是一块货真价实的“细菌社区”，几十上百种菌在里面分层定居，像一座立体的城市。

不过大多数时候，你确实感觉不到它们。没有瘙痒，没有声音。直到近二三十年，科学家不再依赖“把细菌养出来再数”的老办法（很多肠道菌离开肠道就养不活），而是直接用 DNA 测序去“点名”——从样本里读出所有微生物的基因片段，再比对出谁在里面——才真正数清楚这个世界有多大。

\section{数一数：你身上住着多少“房客”}

公布开场悬念的第一部分答案。科学家估计，一个普通成年人肠道里的细菌细胞总数大约是 $3.8\times 10^{13}$ 个——三十八万亿。而你自己身体的人类细胞，大约 $3\times 10^{13}$ 个。两者数量级相同，细菌甚至还略多一点。

先别急着震惊，做个诚实的补充：细菌细胞比人体细胞小得多。一个大肠杆菌的体积大约只有普通人体细胞的千分之一。所以论\textbf{重量}，这些微生物加起来大约只有 $0.2$ 千克——一小袋盐的重量，不是几千克。你猜对了吗？

\begin{qiaoliang}
这个估算怎样做出来？思路你完全可以复现。科学家测出每克粪便大约含 $10^{11}$ 个细菌细胞，再估算大肠内容物总重约几百克，乘一乘就得到肠道细菌总量（身体其他部位的细菌相比之下可以忽略）。类似的估算还可以用在酸奶上：合格的活菌酸奶，每毫升大约含 $10^{8}$ 个以上的活菌。一杯 $150$ 毫升的酸奶，你一口喝下去的活菌就有几百亿个——比全世界人口还多。下次喝酸奶时，可以想想这个数字。
\end{qiaoliang}

数量还不是最惊人的。真正惊人的是\textbf{种类}：你的肠道里大约生活着几百到上千种不同的微生物，它们携带的基因总数，是你自己基因组基因数量的上百倍。从基因的角度看，“你”这座图书馆里，绝大部分藏书根本不是用你的文字写的。这个细节先记住，章末它会成为通往第三册的大门。

\section{从房客到生态系统}

把镜头推进到粒子层——这一章的“粒子”，是一个个微生物细胞，以及它们之间流来流去的分子。

第 57 章讲过，一个细胞就是一座会维护自己化学城市。那么几万亿个细胞挤在你的肠道里，会发生什么？会发生任何拥挤世界里都会发生的事：抢资源、占地盘、合作、互相利用。换句话说，会发生\textbf{生态}。

生态学家研究森林，会区分物种之间的关系。这套词汇搬到你肚子里，几乎一字不改：

\begin{itemize}[itemsep=2pt]
\item \textbf{互利共生}：你吃下自己消化不了的膳食纤维，肠道细菌替你发酵它们，产生的短链脂肪酸反过来养活你的肠壁细胞。你出房租（食物和住所），它交租金（营养和保护）。
\item \textbf{共栖}：很多皮肤细菌住在你身上，吃你的油脂和死皮，对你既无明显好处也无明显坏处——只是蹭住的房客。
\item \textbf{竞争}：同一块地盘、同一种养分，谁先占住就是谁的。肠道里密密麻麻的原住民占满了生态位，外来的坏菌常常挤不进来——健康的菌群本身就是一道防线，是第 61 章免疫防线之外的“民事治安”。
\item \textbf{寄生与致病}：少数住客会翻脸。幽门螺杆菌在胃里钻洞，白色念珠菌在你抵抗力下降时过度繁殖。同一种微生物，平时可能是安静的房客，条件一变就成了捣乱的租客——这叫“条件致病”。
\end{itemize}

所以，你肠道里的世界和一片森林、一个珊瑚礁遵循同一类规则：谁多谁少，取决于食物、空间、酸碱度，以及谁和谁抢、谁和谁合作。你的身体不是一个被细菌“入侵”的纯净城堡，它从一开始就是一个生态系统——而且是一个你和微生物共同演化了亿万年的生态系统。

\section{肠道：一座发酵工厂}

生态系统里最重要的规则，是物质和能量怎么流动。看看你肠道的“产业结构”。

你的小肠负责消化和吸收：淀粉被酶拆成葡萄糖（第 46 章），蛋白质拆成氨基酸，脂肪拆成脂肪酸——这些你在第 55 章的代谢网络里已经见过。但有一类食物你的酶拆不动：膳食纤维，比如豆类、全谷物、蔬菜里的某些复杂碳水化合物。它们穿过小肠，原封不动地进入大肠。

你的酶拆不动，不等于没有酶拆得动。肠道细菌拥有你没有的化学工具箱，能把这些纤维发酵，生成一类叫\textbf{短链脂肪酸}的小分子：乙酸、丙酸、丁酸。别小看这些“细菌尾气”：丁酸是你大肠内壁细胞的主要粮食；乙酸和丙酸进入血液，参与全身的能量代谢和信号传递（第 58 章讲过细胞怎样听消息——细菌的产物也是消息的一种）。

用符号把这件事写下来。酸奶里的乳酸发酵，是细菌版的糖酵解（第 49 章）：一个葡萄糖分子被拆开，净得两个 ATP，剩下的骨架变成乳酸——
\[\chem{C_6H_{12}O_6} \rightarrow 2\,\chem{C_3H_6O_3}\]
注意这个式子两边的原子数是对平的，这就是第 26 章讲过的质量守恒：细菌并没有“消灭”糖，只是把糖里的能量取走一小部分，把剩下的骨架换了个样子排出来。你喝到的酸味，就是这个式子右端的产物。

肠道里的发酵还会生产别的“副产品”。一类是气体：氢气、二氧化碳，有些人的菌群还会制造甲烷——它们就是肠道气体的主要来源，你一天排出的大部分气体不是“吞进去的空气”，而是细菌发酵的烟囱排放。另一类是维生素：部分肠道细菌能合成维生素 K 和好几种 B 族维生素，其中一些能被人体吸收利用——你的“房客”会自己交一部分伙食费。

\begin{center}
\begin{tikzpicture}[scale=1]
\draw[rounded corners, fill=blue!8] (0,0) rectangle (2.6,1.1);
\node[align=center] at (1.3,0.55) {\small 膳食纤维\\ \small（你的酶拆不动）};
\draw[->, thick] (2.6,0.55) -- (3.6,0.55);
\draw[rounded corners, fill=green!12] (3.6,0) rectangle (6.2,1.1);
\node[align=center] at (4.9,0.55) {\small 肠道细菌发酵};
\draw[->, thick] (6.2,0.55) -- (7.2,0.55);
\draw[rounded corners, fill=orange!15] (7.2,0) rectangle (9.8,1.1);
\node[align=center] at (8.5,0.55) {\small 短链脂肪酸};
\draw[->, thick] (8.5,0) -- (8.5,-0.9);
\draw[rounded corners, fill=red!8] (7.2,-2.0) rectangle (9.8,-0.9);
\node[align=center] at (8.5,-1.45) {\small 肠壁细胞的粮食\\ \small 全身的信号分子};
\node[align=center] at (4.9,-2.6) {\small 示意图：颜色和方框只是表示法，不是真实外观};
\end{tikzpicture}
\end{center}

\begin{shiyan}
\textbf{在家做一杯酸奶，亲眼看见细菌工作。}

材料：纯牛奶（巴氏杀菌或常温奶）$500$ 毫升、市售原味活菌酸奶两勺（配料表里有“活菌”字样的）、干净的玻璃罐、保鲜膜。

步骤：① 玻璃罐用开水烫过晾干；② 牛奶加热到约 $40\,^{\circ}\mathrm{C}$（手感温热不烫），拌入两勺酸奶；③ 倒入罐中盖好，包上毛巾放在温暖处（约 $35$～$40\,^{\circ}\mathrm{C}$）静置 $6$～$10$ 小时；④ 凝固后放冰箱冷藏。

观察点：每隔两小时看一眼，记录牛奶从液体变成半固体的过程；闻一闻酸味的变化；想一想：为什么接入“菌种”才行，光把牛奶放温暖处为什么只会坏掉？

安全提示：全程使用干净器具，做好后冷藏并在一两天内吃完；如果闻到腐臭味、看到霉斑，整罐扔掉不要尝。不要用生牛奶做。
\end{shiyan}

\section{一管培养液掀翻一个医学共识}

\begin{zhengju}
你身上的微生物世界里，最出名的“恶棍”大概是幽门螺杆菌——它住在胃里，和胃炎、胃溃疡、胃癌都有关。但在 1982 年以前，医学界几乎众口一词：胃里有强酸（pH 可以低到 1.5），没有细菌能在那里长期生活；胃溃疡的原因是压力太大、吃得太辣、胃酸分泌过多。医生开出的药方是：减压、吃流食、吃抑酸药，严重的做手术切掉一部分胃。

1979 年，澳大利亚病理学家沃伦在胃炎病人的胃黏膜切片里，反复看到一种弯曲的小细菌。他发现这种细菌出现的地方，黏膜几乎总是发炎的。他的同事、年轻医生马歇尔动了真格：从病人胃里把这种细菌培养出来，再去查文献——结果发现“胃是无菌的”这个说法，本身就缺乏过硬的证据。

接下来的故事带有传奇色彩，但请注意它符合科学方法的每一步。马歇尔先尝试用培养出的细菌做动物实验，没能成功致病。1984 年，他做了一个如今绝不会被伦理委员会批准的实验：亲自喝下一管幽门螺杆菌培养液。几天后他开始恶心、呕吐，胃镜检查证实他得了急性胃炎——细菌先、炎症后，时间顺序对上了。更关键的证据是：用抗生素杀灭这种细菌后，许多顽固的胃溃疡患者痊愈，复发率大幅下降。如果溃疡真是“压力造成的”，抗生素凭什么治好它？

替代解释被逐一排除，共识被迫改写。2005 年，沃伦和马歇尔获得诺贝尔奖。这个故事还留下一个值得回味的尾巴：幽门螺杆菌与人类共存了至少几万年，如今科学家发现它可能也有某些“房客”式的另一面（比如与较低的食管反流病发生率相关）——“恶棍”与“住客”的界限，比当年任何人想的都模糊。生态系统里，没有绝对的好菌和坏菌，只有在特定条件下扮演特定角色的成员。
\end{zhengju}

\section{抗生素：一场被按了快进键的选择}

1928 年弗莱明发现青霉素之后，人类第一次握住了对细菌“精确轰炸”的武器。抗生素的原理是化学的：瞄准细菌特有的结构或机器——比如抑制细菌细胞壁的合成，或者卡住细菌核糖体——而这些部件人体细胞没有，所以药物能相对选择性地杀伤细菌。

但请注意“相对”两个字。抗生素进入人体后，分不清肠道里的“好房客”和“坏租客”。一次广谱抗生素疗程，会把肠道菌群的大片街区夷为平地，有时需要几个月才能恢复原来的面貌，有些成员可能永远回不来。这就是为什么医生会嘱咐：抗生素要按疗程、按指征用，不要自己随便吃，也不要吃两天感觉好了就擅自停。

更严重的问题在后面。假设一堆细菌里有极其少数的个体，因为基因的随机差异，碰巧不那么怕某种抗生素。药物一来，敏感的细菌大批死亡，这几个幸运儿活下来，面前突然空出了一大片地盘和食物——它们迅速繁殖，下一代整个群体都不怕这种药了。这就是\textbf{耐药性}，它本质上是自然选择在几星期内上演的实况转播。

这里有一个必须掰清楚的观念：细菌\textbf{不是}“遇到了药，于是努力学会抵抗它”。变异发生在接触药物之前，是随机产生的；药物只是“筛选员”，把碰巧有抵抗力的个体留下来。个体不因需要而改变，改变的是群体里各类型的比例——这条规则，第三册讲进化时还会反复用到。如今，滥用抗生素（包括畜牧业里的滥用）正在把越来越多的细菌筛选成“超级细菌”，这被认为是全球公共卫生的重大威胁之一。

抗生素还会以另一种方式暴露“生态系统”的本质。艰难梭菌是少数能形成耐受孢子的肠道细菌，平时被密密麻麻的原住民压制得抬不起头。一次广谱抗生素疗程把原住民大批清除后，它可能趁机暴发，引起严重的腹泻甚至危及生命——这不是“入侵了一个更强的敌人”，而是“治安队被打垮后，原本不起眼的小混混占领了街区”。治疗这种感染最有效的办法之一，竟是把健康供体的肠道菌群移植回去，让完整的生态系统重建秩序：用药的思路从“杀死细菌”变成了“恢复生态”，这是医学观念的一次深刻转向。

\section{菌群会变：饮食、年龄、药物和环境}

一个生态系统不会静止。你的菌群每天都在随你的生活改变：

\begin{itemize}[itemsep=2pt]
\item \textbf{饮食}：连续几天大量吃肉和连续几天大量吃全谷物、豆类，肠道里的优势菌群会明显不同——发酵纤维的细菌在“高纤维伙食”下兴旺。长期饮食模式的影响更大。
\item \textbf{年龄}：婴儿通过产道和母乳获得第一批微生物“种子”，菌群要花几年才长成接近成人的样子；老年人的菌群多样性通常又会下降。
\item \textbf{药物}：抗生素的影响最直接；一些常见药物（如抑酸药）也会改变肠道环境，间接改变谁能在里面生存。
\item \textbf{环境}：和你一起生活的人、家里的宠物、你接触的土壤和水，都会和你交换微生物。同一屋檐下的家人，菌群往往比陌生人更相似。
\end{itemize}

读到这里，你可能想问：既然菌群影响这么大，那调节菌群是不是能治病、能减肥、能改善情绪？这正是过去十几年科学界最热、也最需要冷静的领域，我们放到边界层说。

\section{边界：这个模型的力量与极限}

“人是生态系统”这个模型非常有用：它解释了为什么抗生素会拉肚子，为什么膳食纤维有益健康，为什么外来病原体有时难以立足。但它也有清晰的边界，而且这一章恰好碰上科学传播里最泛滥的夸大，必须一条条说清楚。

\textbf{第一，相关不等于因果。}大量研究发现，肥胖者、糖尿病患者、抑郁症患者的肠道菌群与健康人“不同”。但不同不代表菌群是原因——也可能是疾病或饮食改变了菌群，甚至两者互为因果。要把“有关”升级成“导致”，需要像马歇尔那样的实验证据：改变菌群，看结果是否随之改变，并排除其他解释。

\textbf{第二，“某一种菌决定性格”这类说法，目前证据不足。}你可能在短视频或营销文章里见过：“肠道是第二大脑，某种菌让你外向”“想变聪明，补这种益生菌”。真实的研究图景是：在无菌小鼠身上，确实观察到菌群缺失会影响行为；个别小型人体研究发现菌群与情绪之间存在关联。但小鼠不是人，小型研究常常无法被重复，而人类性格受基因、成长环境、社会经历等多重因素影响。从“菌群与情绪有关”跳到“某一种菌决定你的性格”，中间隔着一整条尚未建成的证据链。对这类断言，合格的反应是：有趣，待验证，先别掏钱。

\textbf{第三，微生物组研究还很年轻。}不同实验室测同一个人的菌群，方法差异都可能影响结果；很多结论是统计上的平均趋势，落不到某个具体的人身上。医学上真正靠菌群干预站稳脚跟的疗法屈指可数——粪菌移植治疗复发性艰难梭菌感染是最成功的一个，有随机对照试验支持，但它也有严格的适应证和操作规范，绝不是“换菌群治百病”的证据。

\section*{常见误解}

\begin{wuqu}
“益生菌是保健品，多吃总没坏处；想改善体质，喝益生菌饮料就行。”——这个误解有两层。第一层，大多数益生菌产品里的菌，经过胃酸和胆汁的“关卡”后，能在你肠道里长期定居的很少，喝下去多半只是“过路客”；对健康人来说，所谓益处常常缺乏高质量试验的支持。第二层更隐蔽：把肠道菌群想成一个“缺什么补什么”的简单池塘，忽略了它是上千个物种竞争合作的生态系统——外来一两个物种很难改变整个群落。与其花钱买菌，不如先喂好你已有的菌：多样化的蔬菜、水果、全谷物，就是给肠道生态系统最实在的“补贴”。（真有疾病需要治疗时，请听医生的，而不是广告。）
\end{wuqu}

\section{再看那杯酸奶}

回到开场那杯酸奶。现在你可以用生态系统的眼光重读它的配料表：保加利亚乳杆菌和嗜热链球菌，是两个被人类“选中”的房客，在牛奶这个临时生态系统里分工合作——一个先把蛋白质拆开一点，另一个长得更好；它们接力把乳糖变成乳酸，酸到自己都住不下去了，发酵才停下来。你以为你在喝奶，其实你在接收一个微型生态系统的代谢遗产。

而你喝下去的这些菌，绝大多数只是路过。它们要穿过胃酸的强酸关卡（第 52 章讲过的蛋白质变性环境），穿过胆汁，最后能活着抵达大肠的少之又少；即便抵达，面对已经挤满原住民的生态位，也很难定居。这不是说酸奶不好——它是好食物——而是说，你那 $0.2$ 千克、三十八万亿个住客的生态系统，有自己顽固的惯性。

至于那个更大的问题：你是一个人，还是一个生态系统？答案是——你从来就不是“纯人类”。你的线粒体是十几亿年前住进来的细菌（第 57 章内共生的证据），你的肠道是上千个物种的都市，你的皮肤是另一片大陆。“我”这个词，在生物学的放大镜下，边界是模糊的。

\section{第二册的终点，信息的起点}

这本书的第一册从基本粒子讲到化学反应，第二册从碳原子一路走到现在：碳搭出有机分子，有机分子组成大分子，大分子装配成细胞，细胞聚成身体，身体又和微生物组成生态系统。物质的链条一环扣一环，没有一环需要奇迹。

但你可能注意到，这一路走来，有一个问题被我们反复绕开了：\textbf{信息}。细菌为什么会发酵乳糖？因为它带着编码乳糖酶的基因。耐药菌的后代为什么天生耐药？因为耐药的“写法”被复制了下去。你为什么更像你的父母，而不是像你的同桌？你的菌群为什么和你的家人更相似？

这些问题再也不能靠“结构”和“能量”回答了，它们指向另一条线索：生命用什么符号记录配方，怎样复制这份记录，复制出错时又会发生什么。你的三十八万亿微生物房客，每分裂一次就在复制自己的记录；而你身体里的每一个细胞，都装着一份三十亿个字母写成的说明书。第三册就从这里开始：第 63 章先回到 19 世纪的一座修道院菜园，看孟德尔怎样用豌豆把“遗传”从玄学变成算术——然后我们会一路追到 DNA 的双螺旋，追到今天正在改写基因的技术。

\begin{tiaozhan}
生态学家怎样给“菌群的丰富程度”一个数字？一种常用工具叫香农多样性指数：$H=-\sum p_i \ln p_i$，其中 $p_i$ 是第 $i$ 个物种占整个群落的比例。试着算一算：群落 A 有两个物种，比例各占一半；群落 B 也有两个物种，比例是 $99:1$。A 的 $H\approx 0.69$，B 的 $H\approx 0.06$——同样两个物种，均衡程度不同，多样性天差地别。健康人的肠道菌群像 A，多样性高、结构稳；滥用抗生素后的菌群可能滑向 B，一两个物种独大。研究者还常用多样性下降作为菌群“生态失衡”的粗略信号——但记住本章边界层的话：多样性高低与健康的关系也是统计关联，不能对号入座。跳过本节不影响主线。
\end{tiaozhan}

\section*{练一练}

\begin{sikaoti}
\item 画一张“你的一天”物质流图：早餐里的淀粉、膳食纤维和蛋白质分别在哪里被谁（你的酶还是肠道菌）处理，产物分别去了哪里。用箭头标出“你做不到的、细菌替你做”的那一段。
\item 大肠杆菌细胞约重 $10^{-12}$ 克。若一个人肠道细菌约 $4\times 10^{13}$ 个，估算这些细菌的总质量，并和正文“约 $0.2$ 千克”的说法核对。这个估算有哪些不确定的地方？
\item 用自己的话向同桌解释：为什么“细菌接触抗生素后努力学会了耐药”是错的？请画出“用药前—用药后—繁殖后”三个时间点的菌群组成示意（用圆点代表敏感菌、三角代表耐药菌）。
\item 看到一篇题为《科学家发现：这种肠道菌让人更外向》的报道，你会追问哪三个问题来判断它靠不靠谱？（提示：从样本、因果、可重复性入手。）
\item 为什么“胃酸很强，所以胃里没有细菌”当年听起来很有道理，却被推翻了？从这个故事总结：一个看似合理的“机制推断”，为什么仍然需要实验证据检验？
\end{sikaoti}
